% Approach. Preamble needs \usepackage{tikz} and \usetikzlibrary{positioning,arrows.meta,fit,backgrounds}.
% Symbols the Problem section must define (used here, defined there):
%   family \mathcal{F} of instances I with MDP M_I; requirements \Phi, each \varphi \bowtie \theta;
%   rule list R (first match wins, `any` rules); worst/best case w_\varphi^I(R), b_\varphi^I(R);
%   certification (every w passes); training set T \subseteq \mathcal{F}; margin; library \mathcal{L}.

\section{Approach}
\label{sec:approach}

\begin{figure*}[t]
\centering
\resizebox{\textwidth}{!}{%
\begin{tikzpicture}[
  font=\footnotesize,
  >={Stealth[length=2mm]},
  proc/.style={draw, thick, fill=orange!18, minimum width=2.5cm, minimum height=1.05cm, align=center},
  art/.style={draw, rounded corners=3pt, fill=blue!8, minimum width=2.5cm, minimum height=0.7cm, align=center},
  dec/.style={draw, thick, rounded corners=8pt, fill=green!12, minimum width=2.6cm, minimum height=1.05cm, align=center},
  flow/.style={->, thick},
  lbl/.style={font=\scriptsize, align=center},
  region/.style={draw, dashed, rounded corners=6pt, inner sep=7pt},
  manager/.style={draw, rounded corners=8pt, fill=black!3, inner sep=9pt}
]
% Decisions of the adaptation manager
\node[art, minimum width=2.9cm] (inst) at (-0.4, 0) {Instance $I_t$ of\\ the next episode};
\node[dec] (cur) at (4.1, 0) {Current $R_{\mathit{cur}}$\\ certified on $I_t$?};
\node[art] (lib) at (8.4, 1.9) {Library $\mathcal{L}$};
\node[dec] (alt) at (8.4, 0) {Some $R \in \mathcal{L}$\\ certified on $I_t$?};
\node[art, minimum width=2.9cm] (dep) at (18.4, 0) {Rule list controlling\\ the episode on $I_t$};

\draw[flow] (inst) -- node[lbl, above] {model $M_{I_t}$} (cur);
\draw[flow] (cur) -- node[lbl, above] {no} (alt);
\draw[flow] (lib) -- node[lbl, right] {stored rule lists} (alt);
\draw[flow] (cur.north) |- node[lbl, pos=0.62, above] {yes: retain $R_{\mathit{cur}}$} (18.4, 3.2) -- (dep.north);
\draw[flow] (alt) -- node[lbl, above] {yes: switch to the one with the largest margin} (dep);

% Generation, invoked when no stored rule list is certified
\node[proc] (gen) at (5.7, -2.9) {Rule-list\\ generation (LLM)};
\node[proc] (ver) at (9.5, -2.9) {Verification\\ (PRISM)};
\node[proc] (fb)  at (7.6, -4.8) {Feedback};
\node[proc] (sel) at (13.3, -2.9) {Selection\\ and repair};
\node[lbl, anchor=north] (chains) at (7.6, -5.4) {repeated in $K$ independent chains};
\node[dec] (new) at (16.5, -2.9) {$R^\star$ certified\\ on $I_t$?};
\node[art, minimum width=2.9cm] (fail) at (18.4, -5.0) {Adaptation fails,\\ episode deferred};

\draw[flow] (alt.south) -- ++(0, -0.75) -| node[lbl, pos=0.25, above] {no: generate on $T \cup \{I_t\}$} (gen.north);
\draw[flow] (gen) -- node[lbl, above] {rule list $R$} (ver);
\draw[flow] (ver.south) |- node[lbl, pos=0.3, right] {worst-case values,\\ failure locations} (fb.east);
\draw[flow] (fb.west) -| node[lbl, pos=0.7, left] {revise\\ prompt} (gen.south);
\draw[flow] (ver) -- node[lbl, above] {candidates,\\ all rounds} (sel);
\draw[flow] (sel) -- node[lbl, above] {$R^\star$} (new);
\draw[flow] (new.east) -| node[lbl, pos=0.75, right] {yes: deploy,\\ add to $\mathcal{L}$} (dep.south);
\draw[flow] (new.south) |- node[lbl, pos=0.3, left] {no} (fail.west);

\begin{scope}[on background layer]
\node[region, fit=(gen)(ver)(fb)(sel)(chains), label={[lbl, anchor=north east]south east:Generation}] (genbox) {};
\node[manager, fit=(cur)(alt)(lib)(new)(genbox), label={[lbl, anchor=south east]north east:Adaptation manager}] {};
\end{scope}

% Legend
\node[art, minimum width=1.2cm, minimum height=0.4cm] (la) at (-0.6, -3.4) {};
\node[lbl, anchor=west] at (la.east) {artifact};
\node[proc, minimum width=1.2cm, minimum height=0.4cm] (lp) at (-0.6, -4.1) {};
\node[lbl, anchor=west] at (lp.east) {process};
\node[dec, minimum width=1.2cm, minimum height=0.4cm, rounded corners=4pt] (ld) at (-0.6, -4.8) {};
\node[lbl, anchor=west] at (ld.east) {decision};
\end{tikzpicture}%
}
\caption{The adaptation manager deciding which rule list controls the next episode. It retains the current rule list if it is certified on the episode's instance, switches to another certified rule list from its library otherwise, and invokes generation if none is certified. A newly generated rule list is deployed and added to the library only if it is certified on the instance, and otherwise adaptation fails and the episode is deferred.}
\label{fig:overview}
\end{figure*}

\begin{figure}[t]
\begin{algorithmic}[1]
\REQUIRE instance $I_t$ of the next episode, current rule list $R_{\mathit{cur}}$, library $\mathcal{L}$
\IF{$R_{\mathit{cur}}$ is certified on $I_t$}
  \RETURN $R_{\mathit{cur}}$
\ENDIF
\STATE $C \gets \{R \in \mathcal{L} \mid R \text{ is certified on } I_t\}$
\IF{$C \neq \emptyset$}
  \RETURN the rule list in $C$ with the largest smallest margin on $I_t$
\ENDIF
\STATE $T \gets \{I_t\}\, \cup$ the training instances of the $R \in \mathcal{L}$ closest to being certified on $I_t$
\STATE $R^\star \gets \textsc{Generate}(T)$
\IF{$R^\star$ is certified on $I_t$}
  \STATE $\mathcal{L} \gets \mathcal{L} \cup \{R^\star\}$
  \RETURN $R^\star$
\ENDIF
\RETURN \textsc{Defer}$(I_t)$
\end{algorithmic}
\caption{The adaptation policy, run before each episode.}
\label{fig:policy}
\end{figure}

\autoref{fig:overview} shows how the adaptation manager decides which rule list controls the next episode, which takes place on an instance $I_t$ of the family, e.g., a map the system has not seen.
The manager first checks whether the rule list it currently runs is certified on $I_t$ and, if so, retains it.
Otherwise, it checks the rule lists in its library $\mathcal{L}$ and switches to a certified one (\autoref{sec:manager}).
If none is certified, the manager invokes \textit{generation} on a training set that includes $I_t$.
Generation runs a loop in which an LLM writes a rule list (\autoref{sec:generation}), PRISM verifies it on every training instance (\autoref{sec:verification}), and feedback tells the LLM where and why its worst case fails (\autoref{sec:feedback}).
Several independent chains of this loop run on the same training set, and selection and repair picks the best rule list among all their rounds (\autoref{sec:selection}).
If that rule list is certified on $I_t$, it controls the episode and joins the library.
If it is not, adaptation fails and the episode is deferred.
Before deployment, the same generation seeds the library from training sets of small instances, which are cheap to verify.
In terms of a MAPE-K loop~\cite{TODO-Kephart2003-MAPE-K}, the manager's checks form the analysis step, the choice among retaining, switching and generating forms the planning step, and the library is part of the knowledge.
The domain enters this process only through its PRISM models, its requirements and its vocabulary, so the framework itself is domain-independent.

\subsection{Adaptation Manager}
\label{sec:manager}
\label{sec:certification}

\textbf{When Adaptation Happens:}
The system runs \textit{episodes}, e.g., one delivery on one map, and each episode takes place on one instance whose model stays fixed during the episode.
Certificates hold from an instance's initial state, so the manager adapts between episodes, before the next one starts.
Switching controllers during an episode would require certification from the current state and the history that the temporal requirements depend on, together with assumptions about the switch itself, which we leave out of scope.

\textbf{Certification:}
Before the episode on $I_t$, the manager builds the instance's model $M_{I_t}$ from the observed parameters, e.g., the map's size and layout.
For a rule list $R$, it compiles $R$ into a PRISM module that restricts $M_{I_t}$ to the actions the rules allow (\autoref{sec:problem}), and PRISM~\cite{TODO-PRISM} computes the worst case $w_\varphi^{I_t}(R)$ of every requirement $\varphi \in \Phi$ over all behaviors the rules allow.
Certification uses interval iteration~\cite{TODO-Haddad2018-interval-iteration}, which bounds each worst case from both sides.
$R$ is certified on $I_t$ only if the bound on the unfavorable side of every worst case meets its threshold.
A bound interval that contains the threshold makes the result inconclusive, and the manager treats it as not certified.

\textbf{Adaptation Policy:}
\autoref{fig:policy} gives the policy the manager runs before each episode.
It first checks the rule list $R_{\mathit{cur}}$ it currently runs and retains it if it is certified on $I_t$, which avoids switching and costs a single certification (lines~1--3).
Otherwise, it certifies the other rule lists in its library and selects the certified one with the largest smallest margin (lines~4--7), i.e., the one whose worst cases lie furthest beyond their thresholds.
If no rule list is certified, it invokes generation on $I_t$ together with the training instances of the rule list that came closest to being certified on $I_t$, so that the new rule list keeps what that one already handles (lines~8--9).

\textbf{Unsuccessful Adaptation:}
Generation runs a bounded number of chains and rounds, so it can fail even where a suitable rule list exists.
The manager deploys the generated rule list $R^\star$ only if it is certified on $I_t$, and then adds it to the library so that later instances may reuse it (lines~10--13).
Otherwise, adaptation fails, and the manager defers the episode on $I_t$ and records the failure (line~14).
We do not assume a safe fallback controller, since claiming one would require its own model and assumptions.
We describe the extension of the library as a proposed mechanism, which our evaluation exercises only in its sequential study\todo{run the sequential study}.

\textbf{Running Example:}
Consider a library with one rule list per grid-world family, among them the rule list of \autoref{fig:intro-rules}, written for the rooms family.
When the next episode takes place on an 11$\times$11 map of the adjacent family, the rooms rule list reaches the second goal before the third with probability 0.793 in the worst case, just short of the threshold of 0.8, and neither of the other two rule lists is certified either\todo{library of one core run (port\_fam\_history, run 1), target 142}.
The manager therefore invokes generation for this map\todo{report the outcome of generation here once the sequential study has run it}.

\subsection{Rule-List Generation}
\label{sec:generation}

Generation takes a training set $T$ of instances and writes one rule list for all of them.
Before deployment, $T$ is a set of small instances, and at run time, the manager chooses $T$ as described in \autoref{sec:manager}\todo{design choice, not evaluated; the loop supports adding a failing instance to a training set (run.pool\_sets)}.

The LLM receives a prompt that states the problem in natural language and asks for one rule list that works, unchanged, on every training instance.
The prompt describes the family's dynamics, each training instance with its constants and a textual map where the domain provides one, the vocabulary the conditions may use, the actions, and the requirements with their thresholds.
It then explains how a rule list is read, i.e., that the first rule that applies decides and that a rule with action \texttt{any} allows every action for which its condition holds.

Two instructions in the prompt shape what the LLM writes.
The first asks for a \textit{general} strategy, since the rule list will be frozen and used on instances the LLM is not shown, including larger ones.
Conditions may therefore name instance values only through the vocabulary's constants and features, never through numbers that single out one instance (\autoref{sec:problem}).
The second explains that \textit{the worst case decides}.
Wherever no rule applies, or a rule allows several actions, an adversary chooses among the allowed actions, so every reachable state must be covered and every allowed action must be safe.
We found that LLMs otherwise tend to leave states uncovered or to allow actions that stall the agent, which the worst case then exploits\todo{cite or point to an ablation (rq2\_prompting); shown on grid worlds, assumed for the other domains}.
The answer is returned as structured output whose actions are limited to the domain's actions and \texttt{any}~\cite{TODO-structured-outputs}.
\autoref{fig:intro-rules} shows a rule list written in this way while seeding the library for the rooms family.

\subsection{Verification}
\label{sec:verification}

Each training instance gives the rule list its own meaning, since constants and features are evaluated on that instance.
For each $I \in T$, PRISM checks the rule list as in \autoref{sec:certification} and computes both its worst case $w_\varphi^I(R)$ and its best case $b_\varphi^I(R)$ for every requirement.
The best case tells the loop whether some way of completing the rules could still pass, which the feedback reports alongside the worst case.
Within the loop, these values only rank candidates, so PRISM computes them with a fast iterative method, and the selected rule list is certified with the bounded method before it is deployed.

\subsection{Feedback}
\label{sec:feedback}

When the rule list passes on every training instance with a margin to spare, or after a fixed number of rounds, the chain stops.
Otherwise, the LLM receives a revise prompt with the worst-case results on each training instance and, for each failing requirement, the states where the worst-case run is lost.
A state is \textit{lost} for a requirement when no behavior can meet the requirement from there any more, e.g., after a collision.
For each state from which much of the worst case's probability flows into lost states, the feedback names the rule that decided the state, the actions it allowed, the action the adversary took, and for every action the features that hold for it and its risk of reaching a lost state in one step.
The prompt also lists the earlier attempts of the chain with their results, and the LLM returns a complete revised rule list.

\begin{figure}[t]
\centering
\fbox{\begin{minipage}{0.95\columnwidth}
\footnotesize\ttfamily\raggedright
Instance 1: FAILS avoid\_moving\_seg1\\
- avoid\_moving\_seg1 = 0.187 (needs >= 0.7).\\
\hspace*{1em}* 0.193 of runs fail from here: agent at (0,1), g1=false, g2=false, g3=false, obs\_idx=3 (obstacle at (1,2), next at (0,2)); decided by rule 3 `!g1 \& !wall \& !obstacle \& !early \& closer -> any'; allowed down; the worst case takes down.\\
\hspace*{2em}actions here: up: [wall, risky] risk: collision 0.15; right: [toward, closer, patrol, obstacle] risk: collision 0.70; down: [risky, toward, closer] risk: collision 0.15; left: [no feature true]
\end{minipage}}
\caption{An excerpt of the feedback on an earlier rule list while seeding the rule list of \autoref{fig:intro-rules}, for the requirement to avoid the moving obstacle until the first goal. In its next answer, the LLM added \texttt{!risky} to the rule that decided this state.}
\label{fig:feedback}
\end{figure}

\autoref{fig:feedback} shows an excerpt of this feedback from seeding the rule list of \autoref{fig:intro-rules}.
On the smallest training instance, an earlier rule list avoided the moving obstacle until the first goal with probability 0.187 in the worst case, against a threshold of 0.7.
From the state shown, the rule deciding there sends the agent \texttt{down}, a move whose slips can collide with the obstacle, as its \texttt{risky} feature and its one-step risk show.
In its next answer, the LLM added \texttt{!risky} to that rule, which removes such moves wherever the rule applies on any instance.
The feedback never shows which action is optimal.
It locates the failure and describes the actions in the vocabulary's terms, so that the LLM has to repair its strategy rather than copy single decisions, which would not carry over to other instances.

\subsection{Selection and Repair}
\label{sec:selection}

A single chain may stall on a poor first answer, so we run $K$ independent chains of the loop on each training set.
Once they end, we collect the best distinct rule lists from all rounds of all chains as candidates.
Each candidate is then \textit{repaired} by hill climbing over small edits, i.e., dropping a rule, adding or dropping a conjunct, changing an action, or swapping two neighboring rules.
An edit is kept when it increases the rule list's margins on the training instances.
Repair scores edits with a fast exact checker of the training instances, so that it can try many edits without running PRISM.
Among the repaired candidates, we select the rule list that passes on the most training instances and, among those, the one with the largest smallest margin.
Passing on every training instance matters, since such rule lists carry over to many more of the instances the system meets later than those that do not\todo{shown with the previous loop on the grid-world families (14.2 vs.\ 5.2 per list), Pac-Man (4.0 vs.\ 1.8) and Mars Exploration Rovers across $n$ (6 vs.\ 0); assumed for random grids; re-measure with the core loop}.

Where an \texttt{any} rule allows several actions, a deployed controller needs to choose one.
Selection fixes one order of preference among the actions, the one under which the rule list does best on the training instances, and the frozen rule list takes the first allowed action in that order\todo{REVIEW: the verifier choosing the order; the evaluation also reports certification with the choice left open}.
The certificate of the frozen rule list therefore covers the controller with this order, and not every way of resolving the \texttt{any} rules of the rule list the LLM wrote.
When the manager invoked generation for $I_t$, it certifies the selected rule list on $I_t$ before deploying it (\autoref{sec:manager}).
\todo{the evaluation approximates this loop: it seeds a library and certifies it on held-out instances, without the run-time extension; say so in the evaluation setup}
